% GENERATED FILE. Edit canonical lessons, then run npm run generate:books.
% canonical-source-hash: fnv1a64:70b4f82c9d7e968a
% canonical-lessons: GU-C292-thag, GU-C292-kaubhad, GU-C292-lac, GU-C292-jasus, GU-C292-jhadapvu, GU-R292-review-285-292

\chapter{A Swindle Uncovered}
\label{ch:gu-a-swindle-uncovered}
\hlchaptermodality{292}

\begin{chapteropening}
\textbf{By the end of this chapter:} \emph{I can tell how a swindler was caught.}

Complete the last lesson of chapter 292: I can tell how a swindler was caught.
\end{chapteropening}

\section[\texorpdfstring{\gu{ઠગ} (ṭhag)}{ઠગ (ṭhag)}]{\gu{ઠગ} (ṭhag) — a swindler}

\label{lesson:GU-C292-thag}



\begin{quote}
Before the new one: say the Gujarati for a prisoner, then the Gujarati for to set free.
\end{quote}

\subsection*{You'll want to know: \gu{ઠગ}}
\textbf{\gu{ઠગ}} — \emph{ṭhag} — \textquotedblleft{}a swindler\textquotedblright{}.

\begin{quote}
\gu{એક} \gu{ઠગે} \gu{ગામના} \gu{લોકોને} \gu{નકલી} \gu{નોટો} \gu{આપી}.
\end{quote}

\begin{quote}
\emph{ek ṭhage gāmnā lokone naklī noṭo āpī.}
\end{quote}

\begin{quote}
A swindler gave the villagers fake banknotes.
\end{quote}

A masculine or feminine noun, by the person meant.

\begin{cousinweb}
\textbf{\gu{ઠગ}} (\emph{ṭhag}) comes from \emph{*ṭhagg}, \textquotedblleft{}to cheat\textquotedblright{}, a form CDIAL reconstructs, and CDIAL lists the Gujarati noun (CDIAL 5489). English \emph{thug} was taken from this family of words.
\end{cousinweb}

\begin{grammarlens}[title={What you've built}]
One more word for telling what happened.
\end{grammarlens}

\subsection*{Your turn}
\begin{itemize}
\raggedright
  \item \emph{Say it:} \emph{ṭhag}
  \item \emph{Say it:} \emph{ṭhag}, once more
  \item \emph{Say it:} \emph{ek ṭhage gāmnā lokone naklī noṭo āpī.} — A swindler gave the villagers fake banknotes.
  \item \emph{Recall:} say the Gujarati for a prisoner, then the Gujarati for to set free, then say \emph{ṭhag} again
\end{itemize}

\begin{tcolorbox}[breakable,colback=teal!4,colframe=teal!35!black,title={Before you move on}]
What does \gu{ઠગ} mean? (\textquotedblleft{}A swindler\textquotedblright{}.) Say it once more.
\end{tcolorbox}
\section[\texorpdfstring{\gu{કૌભાંડ} (kaubhā̃ḍ)}{કૌભાંડ (kaubhā̃ḍ)}]{\gu{કૌભાંડ} (kaubhā̃ḍ) — a scam}

\label{lesson:GU-C292-kaubhad}



\begin{quote}
Before the new one: say the Gujarati for to set free, then the Gujarati for a swindler.
\end{quote}

\subsection*{You'll want to know: \gu{કૌભાંડ}}
\textbf{\gu{કૌભાંડ}} — \emph{kaubhā̃ḍ} — \textquotedblleft{}a scam\textquotedblright{}.

\begin{quote}
\gu{આ} \gu{કૌભાંડમાં} \gu{બે} \gu{મોટા} \gu{વેપારી} \gu{હતા}.
\end{quote}

\begin{quote}
\emph{ā kaubhā̃ḍmā̃ be moṭā vepārī hatā.}
\end{quote}

\begin{quote}
Two big merchants were part of this scam.
\end{quote}

A neuter noun.

\begin{cousinweb}
\textbf{\gu{કૌભાંડ}} (\emph{kaubhā̃ḍ}) No dictionary we use records an older form, so its history is not settled. It is a large fraud, usually one with money and people in high places in it.
\end{cousinweb}

\begin{grammarlens}[title={What you've built}]
One more word for telling what happened.
\end{grammarlens}

\subsection*{Your turn}
\begin{itemize}
\raggedright
  \item \emph{Say it:} \emph{kaubhā̃ḍ}
  \item \emph{Say it:} \emph{kaubhā̃ḍ}, once more
  \item \emph{Say it:} \emph{ā kaubhā̃ḍmā̃ be moṭā vepārī hatā.} — Two big merchants were part of this scam.
  \item \emph{Recall:} say the Gujarati for to set free, then the Gujarati for a swindler, then say \emph{kaubhā̃ḍ} again
\end{itemize}

\begin{tcolorbox}[breakable,colback=teal!4,colframe=teal!35!black,title={Before you move on}]
What does \gu{કૌભાંડ} mean? (\textquotedblleft{}A scam\textquotedblright{}.) Say it once more.
\end{tcolorbox}
\section[\texorpdfstring{\gu{લાંચ} (lā̃c)}{લાંચ (lā̃c)}]{\gu{લાંચ} (lā̃c) — a bribe}

\label{lesson:GU-C292-lac}



\begin{quote}
Before the new one: say the Gujarati for a swindler, then the Gujarati for a scam.
\end{quote}

\subsection*{You'll want to know: \gu{લાંચ}}
\textbf{\gu{લાંચ}} — \emph{lā̃c} — \textquotedblleft{}a bribe\textquotedblright{}.

\begin{quote}
\gu{એક} \gu{સિપાઈએ} \gu{ઠગ} \gu{પાસેથી} \gu{લાંચ} \gu{લીધી}.
\end{quote}

\begin{quote}
\emph{ek sipāīe ṭhag pāsethī lā̃c līdhī.}
\end{quote}

\begin{quote}
A constable took a bribe from the swindler.
\end{quote}

A feminine noun.

\begin{cousinweb}
\textbf{\gu{લાંચ}} (\emph{lā̃c}) goes back to Sanskrit \emph{lañcā}, \textquotedblleft{}a present, a bribe\textquotedblright{}, and CDIAL lists the Gujarati noun (CDIAL 10914). It is feminine: \textbf{\gu{લાંચ} \gu{લીધી}}, took a bribe.
\end{cousinweb}

\begin{grammarlens}[title={What you've built}]
One more word for telling what happened.
\end{grammarlens}

\subsection*{Your turn}
\begin{itemize}
\raggedright
  \item \emph{Say it:} \emph{lā̃c}
  \item \emph{Say it:} \emph{lā̃c}, once more
  \item \emph{Say it:} \emph{ek sipāīe ṭhag pāsethī lā̃c līdhī.} — A constable took a bribe from the swindler.
  \item \emph{Recall:} say the Gujarati for a swindler, then the Gujarati for a scam, then say \emph{lā̃c} again
\end{itemize}

\begin{tcolorbox}[breakable,colback=teal!4,colframe=teal!35!black,title={Before you move on}]
What does \gu{લાંચ} mean? (\textquotedblleft{}A bribe\textquotedblright{}.) Say it once more.
\end{tcolorbox}
\section[\texorpdfstring{\gu{જાસૂસ} (jāsūs)}{જાસૂસ (jāsūs)}]{\gu{જાસૂસ} (jāsūs) — a detective; a spy}

\label{lesson:GU-C292-jasus}



\begin{quote}
Before the new one: say the Gujarati for a scam, then the Gujarati for a bribe.
\end{quote}

\subsection*{You'll want to know: \gu{જાસૂસ}}
\textbf{\gu{જાસૂસ}} — \emph{jāsūs} — \textquotedblleft{}a detective; a spy\textquotedblright{}.

\begin{quote}
\gu{જાસૂસે} \gu{કહ્યું} \gu{કે} \gu{જો} \gu{ઠગ} \gu{હજી} \gu{શહેરમાં} \gu{છે} \gu{તો} \gu{પોલીસ} \gu{તેને} \gu{પકડશે}.
\end{quote}

\begin{quote}
\emph{jāsūse kahyũ ke jo ṭhag hajī śahermā̃ che to polīs tene pakaḍśe.}
\end{quote}

\begin{quote}
The detective said that if the swindler was still in the city, the police would catch him.
\end{quote}

A masculine or feminine noun, by the person meant.

\begin{cousinweb}
\textbf{\gu{જાસૂસ}} (\emph{jāsūs}) is usually traced, through Persian, to Arabic \emph{jāsūs}, \textquotedblleft{}a spy\textquotedblright{}. In stories it is the detective who works out who did it.
\end{cousinweb}

\begin{grammarlens}[title={What you've built}]
One more word for telling what happened.
\end{grammarlens}

\subsection*{Your turn}
\begin{itemize}
\raggedright
  \item \emph{Say it:} \emph{jāsūs}
  \item \emph{Say it:} \emph{jāsūs}, once more
  \item \emph{Say it:} \emph{jāsūse kahyũ ke jo ṭhag hajī śahermā̃ che to polīs tene pakaḍśe.} — The detective said that if the swindler was still in the city, the police would catch him.
  \item \emph{Recall:} say the Gujarati for a scam, then the Gujarati for a bribe, then say \emph{jāsūs} again
\end{itemize}

\begin{tcolorbox}[breakable,colback=teal!4,colframe=teal!35!black,title={Before you move on}]
What does \gu{જાસૂસ} mean? (\textquotedblleft{}A detective; a spy\textquotedblright{}.) Say it once more.
\end{tcolorbox}
\section[\texorpdfstring{\gu{ઝડપવું} (jhaḍapvũ)}{ઝડપવું (jhaḍapvũ)}]{\gu{ઝડપવું} (jhaḍapvũ) — to nab}

\label{lesson:GU-C292-jhadapvu}



\begin{quote}
Before the new one: say the Gujarati for a bribe, then the Gujarati for a detective.
\end{quote}

\subsection*{You'll want to know: \gu{ઝડપવું}}
\textbf{\gu{ઝડપવું}} — \emph{jhaḍapvũ} — \textquotedblleft{}to nab\textquotedblright{}.

\begin{quote}
\gu{છેવટે} \gu{પોલીસે} \gu{ઠગને} \gu{સ્ટેશન} \gu{પર} \gu{ઝડપ્યો}.
\end{quote}

\begin{quote}
\emph{chevṭe polīse ṭhagne sṭeśan par jhaḍapyo.}
\end{quote}

\begin{quote}
In the end the police nabbed the swindler at the station.
\end{quote}

Past, with \gu{તેણે} (he or she …): \gu{તેણે} \gu{ઠગને} \gu{ઝડપ્યો}, \emph{teṇe ṭhagne jhaḍapyo}.

\begin{cousinweb}
\textbf{\gu{ઝડપવું}} (\emph{jhaḍapvũ}) comes from \emph{*jhaṭ}, a form CDIAL reconstructs for a quick, sudden movement; CDIAL lists Gujarati \emph{jhaṛapvũ}, \textquotedblleft{}to snatch\textquotedblright{} (CDIAL 5327). In news reports it is what the police do to a wanted man: they grab him fast.
\end{cousinweb}

\begin{grammarlens}[title={What you've built}]
One more word for telling what happened.
\end{grammarlens}

\subsection*{Your turn}
\begin{itemize}
\raggedright
  \item \emph{Say it:} \emph{jhaḍapvũ}
  \item \emph{Say it:} \emph{jhaḍapvũ}, once more
  \item \emph{Say it:} \emph{chevṭe polīse ṭhagne sṭeśan par jhaḍapyo.} — In the end the police nabbed the swindler at the station.
  \item \emph{Recall:} say the Gujarati for a bribe, then the Gujarati for a detective, then say \emph{jhaḍapvũ} again
\end{itemize}

\begin{tcolorbox}[breakable,colback=teal!4,colframe=teal!35!black,title={Before you move on}]
What does \gu{ઝડપવું} mean? (\textquotedblleft{}To nab\textquotedblright{}.) Say it once more.
\end{tcolorbox}
\section[(dialogue)]{Looking back: from A Thief in the Night to A Swindle Uncovered}

\label{lesson:GU-R292-review-285-292}



\begin{quote}
Say the Gujarati for a thief and a theft.
\end{quote}

\subsection*{Your turn}
Say these aloud:

\begin{itemize}
\raggedright
  \item a thief, a theft, to flee, to catch, people
  \item a crime, a constable, a suspicion, very, to hide (something)
  \item a trace, a footstep, to knock, to hide oneself, to break (something)
  \item an arrest, to confess, to deny, a criminal, a jail
  \item a court, a lawyer, justice, a punishment, to undergo
  \item a bandit, a gun, a threat, to threaten, to snatch
  \item a kidnapping, missing, to surround, a prisoner, to set free
  \item a swindler, a scam, a bribe, a detective, to nab
\end{itemize}

\begin{tcolorbox}[breakable,colback=teal!4,colframe=teal!35!black,title={Before you move on}]
A thief? (\textbf{\gu{ચોર}}, \emph{cor}.) A theft? (\textbf{\gu{ચોરી}}, \emph{corī}.) To flee? (\textbf{\gu{નાસવું}}, \emph{nāsvũ}.) To catch? (\textbf{\gu{પકડવું}}, \emph{pakaḍvũ}.) People? (\textbf{\gu{લોકો}}, \emph{loko}.) A crime? (\textbf{\gu{ગુનો}}, \emph{guno}.) A constable? (\textbf{\gu{સિપાઈ}}, \emph{sipāī}.) A suspicion? (\textbf{\gu{શંકા}}, \emph{śankā}.) Very? (\textbf{\gu{ખૂબ}}, \emph{khūb}.) To hide (something)? (\textbf{\gu{છુપાવવું}}, \emph{chupāvvũ}.) A trace? (\textbf{\gu{ભાળ}}, \emph{bhāḷ}.) A footstep? (\textbf{\gu{પગલું}}, \emph{paglũ}.) To knock? (\textbf{\gu{ખખડાવવું}}, \emph{khakhḍāvvũ}.) To hide oneself? (\textbf{\gu{સંતાવું}}, \emph{santāvũ}.) To break (something)? (\textbf{\gu{તોડવું}}, \emph{toḍvũ}.) An arrest? (\textbf{\gu{ધરપકડ}}, \emph{dharpakaḍ}.) To confess? (\textbf{\gu{કબૂલવું}}, \emph{kabūlvũ}.) To deny? (\textbf{\gu{નકારવું}}, \emph{nakārvũ}.) A criminal? (\textbf{\gu{ગુનેગાર}}, \emph{gunegār}.) A jail? (\textbf{\gu{જેલ}}, \emph{jel}.) A court? (\textbf{\gu{અદાલત}}, \emph{adālat}.) A lawyer? (\textbf{\gu{વકીલ}}, \emph{vakīl}.) Justice? (\textbf{\gu{ન્યાય}}, \emph{nyāy}.) A punishment? (\textbf{\gu{સજા}}, \emph{sajā}.) To undergo? (\textbf{\gu{ભોગવવું}}, \emph{bhogavvũ}.) A bandit? (\textbf{\gu{ડાકુ}}, \emph{ḍāku}.) A gun? (\textbf{\gu{બંદૂક}}, \emph{bandūk}.) A threat? (\textbf{\gu{ધમકી}}, \emph{dhamkī}.) To threaten? (\textbf{\gu{ધમકાવવું}}, \emph{dhamkāvvũ}.) To snatch? (\textbf{\gu{છીનવવું}}, \emph{chīnavvũ}.) A kidnapping? (\textbf{\gu{અપહરણ}}, \emph{apaharaṇ}.) Missing? (\textbf{\gu{ગુમ}}, \emph{gum}.) To surround? (\textbf{\gu{ઘેરવું}}, \emph{ghervũ}.) A prisoner? (\textbf{\gu{કેદી}}, \emph{kedī}.) To set free? (\textbf{\gu{છોડાવવું}}, \emph{choḍāvvũ}.) A swindler? (\textbf{\gu{ઠગ}}, \emph{ṭhag}.) A scam? (\textbf{\gu{કૌભાંડ}}, \emph{kaubhā̃ḍ}.) A bribe? (\textbf{\gu{લાંચ}}, \emph{lā̃c}.) A detective? (\textbf{\gu{જાસૂસ}}, \emph{jāsūs}.) To nab? (\textbf{\gu{ઝડપવું}}, \emph{jhaḍapvũ}.)
\end{tcolorbox}
